\input format.tex

\usepackage{graphicx}
\graphicspath{{cores/}}

\usepackage{environ}
\usepackage{colortbl,array,booktabs}
\usepackage{tabularx}

\colorlet{TablaBordeSuperior}{topcolor}
\colorlet{TablaBordeInferior}{topcolor}
\colorlet{TablaCentroSuperior}{blue!1}
\colorlet{TablaCentroInferior}{blue!20}
\colorlet{FuenteCabeceraTabla}{white}

\newcolumntype{M}[1]{>{\centering\arraybackslash}m{#1}}

\tcbset{rtab/.style={
freelance,
frame code={
 \path[top color=topcolor,bottom color=topcolor]
   ([yshift=-#1*(\baselineskip+2pt)]interior.north west) --
   ([yshift=-#1*(\baselineskip+2pt)]interior.north east) {[sharp corners]--
    ([yshift=3pt]interior.north east) --
    ([yshift=3pt]interior.north west)} -- cycle;

  },
interior code={},
 }
}

\newcommand\fuentecabecera[1]{\textcolor{black}{\textbf{#1}}}

\NewEnviron{RCtable}[1]
  {%
    \begin{tcolorbox}[#1]\BODY\end{tcolorbox}%
  }



\begin{document}

\vspace*{5mm}
%% 各章节
\setlength{\arrayrulewidth}{1pt}
\fontsize{9.3pt}{11pt}\selectfont
\color{gray2}

%{\noindent\bf\sanhao 二、肠道菌含量}

%\vspace*{6mm}

\begin{spacing}{1.5}
\begin{LRaside}[.8]{2.3 免疫指数}
\noindent
\includegraphics[width=\linewidth]{hanliang.pdf}
\asidebreak %
免疫系统是人体最重要的防线。有七成以上的免疫细胞，如巨噬细胞、T细胞、NK细胞、B细胞等，分布在肠道各处。肠道免疫细胞产生七成以上的免疫球蛋白A，用来保护人体。综上所述，肠道是人体名副其实最重要的免疫器官。\\
肠道微生物与宿主免疫系统之间协同进化，调节B细胞、T细胞等免疫细胞功能。建立和维持肠道微生物与宿主免疫系统之间有益的相互作用是保障机体健康的关键。如果肠道菌活动发生紊乱，它们将变成致病的源头。一旦两者间的平衡被打破，各种免疫炎症性疾病，可能接踵而来。
\end{LRaside}
\end{spacing}

\vspace*{4mm}
%\noindent
%2.1 您的常见肠道菌含量
%\vspace*{1.5mm}

%\begin{RCtable}{
%  rtab=1,
%  tabularx*={\arrayrulecolor{gray2}}%
%    {M{0.3\linewidth} M{0.31\linewidth}M{0.34\linewidth}}
%}
%\fuentecabecera{测试项目} & \fuentecabecera{健康状态提示} & \fuentecabecera{检测结果评价} \\[3pt]
%\bf{测试项目} & \bf{健康状态提示} & \bf{检测结果评价} \\[3pt]
%\end{RCtable}

\begin{tctabularx}{tabularx={m{4.8cm}<{\centering}m{5.2cm}<{\centering}m{4.55cm}<{\centering}}}
&&
\\[-6pt]
\cellcolor{topcolor} \raisebox{2.614pt}{\color{gray0}\bf 检测项目} &
\cellcolor{topcolor} \raisebox{2.614pt}{\color{gray0}\bf 健康状态提示} &
\cellcolor{topcolor} \raisebox{2.614pt}{\color{gray0}\bf 检测结果评价}
\end{tctabularx}

\vspace*{-4.25mm}
\fontsize{8.8pt}{11pt}\selectfont
\arrayrulecolor{gray2}
\begin{longtable}{m{4.8cm}m{5.2cm}<{\centering}m{0cm}@{}m{4.61cm}<{\centering}}
\hline
\parbox[c]{\hsize}{\vskip7pt {\lantxh 免疫指数\\肠道免疫功能} \vskip7pt} & \parbox[c]{\hsize}{\vskip7pt\centerline{\raisebox{-1.5ex}{\includegraphics[width=5cm,height=0.55cm]{bar.pdf}}}\vskip7pt}  &
\hspace*{-3.17cm}\raisebox{-0.45ex}{\includegraphics{cry.pdf}}
 & \begin{minipage}{4.60cm}\begin{center}{{\color{orange}\lantxh 偏低{\\ \bahao 表明免疫力可能降低，增加疾病风险}} }\end{center} \end{minipage} \\
% & \begin{minipage}{4.60cm}\begin{center}{{\lantxh 偏低{\\ \bahao 表明免疫力可能降低，增加疾病风险}} }\end{center} \end{minipage} \\
% & \begin{minipage}{4.60cm}\begin{center}{{\color{red}\lantxh 偏低 \\ \bahao 表明免疫力可能降低，增加疾病风险} }\end{center} \end{minipage} \\
\hline
\multicolumn{3}{l}{}\\
\multicolumn{3}{l}{\includegraphics[scale=1]{greencircle.pdf} {\fontsize{9.3pt}{11pt}\selectfont\bf 绿色表示健康}   \includegraphics[scale=1]{orangecircle.pdf}  {\fontsize{9.3pt}{11pt}\selectfont\bf 橙色表示需要关注}  \includegraphics[scale=1]{redcircle.pdf} {\fontsize{9.3pt}{11pt}\selectfont\bf 红色表示有风险}}
\end{longtable}

\vspace*{6mm}
\begin{spacing}{1.5}
\begin{LRaside}[.8]{结果分析}
\noindent
\includegraphics[width=\linewidth]{result_fenbu.pdf}
\asidebreak %
综合您的检测结果，
您肠道内的免疫指数指标异常，不利于人体和肠道健康。肠道屏障功能可能较弱，对炎症和致病菌等防御力较弱，请注意预防改善。
\end{LRaside}
\end{spacing}

\end{document}

